\documentclass[book.tex]{subfiles}

\begin{document}

\chapter{A Brief Introduction to Superconductor Theory}

\label{chap:superconductors}
\noindent\rule{11cm}{0.4pt} \\
\textbf{Prerequisites:}  \autoref{chap:schrodinger} \\
\textbf{Difficulty Level:} ** \\
\noindent\rule{11cm}{0.4pt}
\vspace{5mm}

In this chapter we provide a brief introduction to the quantum mechanics of superconductors. A superconductor is a special system which is both a perfect conductor and which expels external magnetic fields.


\section{BCS Theory}
The BCS theory provides a basic model of superconductivity by modeling the interactions between electrons and phonons in the system. In particular, phonons mediate an attractive electron-electron interactions. These interacting pairs are called Cooper pairs. This system can be modeled by an attractive potential $V(r_1 - r_2)$ where $r_1$ and $r_2$ are the positions of two electrons.

Superconductivity occurs when this attractive force dominates the repulsive force from the Coulomb potential. The BCS theory also successfully predicts the second order phase transition at the critical temperature $T_c$.

\section{Superconducting Density Functional Theory}

A successful computational strategy to model superconductivity has been to use density functional theory with a custom exchange-correlation functional to model superconducting systems.

%\textcolor{red}{TODO: Say more about custom xc-functionals for superconducting systems.}


\end{document}